\paragraph{What is held aside.}
A Datadog widget is a list of named queries plus a formula over them, and this paper is
about the metrics query language rather than the formula layer. Scoping it out is cheap
and the corpus says so: of $6{,}022$ formulas ($444$ distinct strings), $75.6\%$ are a
bare reference to a single named query, $13.8\%$ apply one function, $13.2\%$ contain
arithmetic and $8.2\%$ a numeric constant. The semantics sit in the queries, not above
them.

The population is the metric formula layer and not the whole of it, which is the
correction the sentence above needed. A formula sits over a request's queries, so it
belongs to the metrics language only when they do, and the extractor applied that filter
to the queries and not to the formulas. The $11{,}628$ formulas an earlier version of
this paragraph reported were $6{,}022$ metric ones and $5{,}606$ over log, event,
security-signal and other requests --- and the second group is $90.1\%$ bare references,
so it diluted every share here towards the answer the paragraph wanted. Arithmetic in the
formula layer is nearly twice what was claimed. The conclusion survives: three quarters of
the metric formula layer is a name, and nothing in it computes what a query cannot.

\paragraph{What the corpus is not.}
It is a census of \emph{published} dashboards, not of production estates. Vendor
integration dashboards are written to demonstrate an integration, and community
dashboards are written to be shared, so both are plausibly tidier and more
template-driven than a dashboard grown inside one company. The direction of that bias is
not obvious --- a tidier corpus is easier to parse but no easier to translate --- and
Section~\ref{sec:threats} returns to it.

